% Kapitel 9 -- Von Kuonrâts Frauenraub vor Znaim

\section{Von Kuonrâts Frauenraub vor Znaim}

\mylettrine{I}{n} denselben Tagen, da man zu \gls{Hardegg} noch flüsternd von der armen \gls{Regina} sprach, und von jenem Reiter, der, wie es hieß, seit seinem Sturz bei Vollmond am Felsen wiederkehre, verschied im Herzogtum \gls{Oesterriche} auch der Babenberger Herzog, den viele den Glorreichen nannten. So fiel auf die Grenzherrschaften an Thaya und Kamp ein doppelter Schatten, aus Schuld und aus Geschichte gewoben. Noch war das \gls{Interregnum} nicht angebrochen, und doch erkennt, wer von den späteren Wirrnissen her zurückblickt, in jenem Jahre schon den dumpfen Vorboten des Unfriedens. Die große Ordnung stand noch, aber ihre Balken knarrten bereits.

Da \gls{Kuonrat von Steinare} nach jenem Unheil aus \gls{Hardegg} hinweg und mit \gls{Gruober} nach Böhmen gesandt worden war, traf ihn die Fortsetzung seines Unheils nicht mehr unter den Zinnen seines Herrn, sondern auf fremdem Boden. Er ritt mit \gls{Gruober} in den böhmischen Landen gegen \gls{Znaim} umher. Gemeinsam führten sie zehn Mann bei sich. Unter ihnen war auch Kuonr\^{a}ts \gls{Knecht}. Schon seit zwei Tagen waren sie zwischen kleinen Weilern, kahlen Obstgärten und dunklen Weingärten unterwegs, und über allem lag jener fremde Laut der Gegend, in welchem dem jungen Herrn jedes Haus wie eine verschlossene Truhe vorkam. Als nun die Nacht vollends niedersank, von dünnem Mondlicht und fahrendem Nebel durchzogen, lagen die einzelnen Höfe weit voneinander, als hätten sie sich selber vor bösen Nachbarn verborgen.

Als sie an einen Hohlweg kamen, der sich zu mehreren Weilern verzweigte, hielt \gls{Gruober} an und teilte seine Leute wie ein Wolfsjäger, der den Kreis um die Beute zieht. \enquote{Dort drüben, unter den Weiden}, sprach er zu \gls{Kuonrat von Steinare}, \enquote{liegt ein Hof, klein genug für zwei Männer. Du nimmst den \gls{Knecht} und siehst zu, dass du mir wenigstens eine Magd herausholst. Mir ist nicht um Schönheit zu tun; jung soll sie sein und kräftig. Wir anderen streichen die übrigen Weiler ab und sammeln uns wieder vor dem ersten Hahnenschrei.} Damit ritt der Alte davon, und seine Leute verschwanden zwischen dunklen Schobern.

Jeder kluge Räuber hätte sich nun geduckt, den Hof ausgespäht und die Tür erst berührt, wenn er gewußt hätte, wer darin wache. \gls{Kuonrat von Steinare} aber hielt entweder offenes Auftreten für eine Art Tapferkeit und Torheit für männliche Kühnheit, oder aber er glaubte wirklich, dass \gls{Gruober} nur ein Geschäftsmann sei und die Frauen freiwillig mit ihm kommen würden, hörten sie nur sein Angebot. Was auch immer er sich nun dachte oder nicht dachte, so ging er nun, statt sich anzuschleichen, geraden Schrittes auf das Haus zu, schlug mit der Faust an das Holz und rief mit lauter Stimme: \enquote{Tut auf! Tut auf! Ich suche Mägde, die Dienst brauchen. Wer eine ledige Magd entbehren kann, der soll sie mir zeigen. Brot und Kleid sollen ihr nicht fehlen.}

Die Leute des Hofes sprachen die slawische Zunge der Gegend, und \gls{Kuonrat von Steinare} verstand sie so wenig, wie sie ihn verstanden. Doch die Gerüchte von verschwindenden Mädchen waren längst bis zu solchen abgelegenen Höfen gedrungen. Ein alter Bauer, der einmal auf dem Markt bei \gls{Znaim} mit deutschen Kaufleuten geredet hatte, vernahm aus Kuonrats Worten nur Bruchstücke, das Wort \enquote{Magd}, das Locken mit Brot, den Ton des Begehrens, und das genügte. Sie wunderten sich wohl, weshalb ein Frauenräuber anklopfte wie ein ehrbarer Werber; eben darum hielten sie es für eine List und beschlossen, ihn hereinzulassen.

Die Tür öffnete sich einen Spalt, und derselbe alte Mann machte mit der Hand eine einladende Bewegung. \gls{Kuonrat von Steinare} trat ohne Argwohn hinein, das Herz voller törichter Zuversicht. Im selben Augenblick warfen zwei Bauern einen groben Mühlsack über seinen Kopf und die Schultern, zogen ihn nieder und wickelten das Tuch um Arme und Hals, dass er wie ein gefesseltes Kalb zu Boden stürzte. Zugleich schnellte draußen von der Seite eines Schobers die Sehne eines Bogens. Der Pfeil fuhr gegen den \gls{Knecht}, traf aber den hochgerissenen Schild, drang ins Holz und blieb zitternd stecken.

Hier zeigte der \gls{Knecht} mehr Mut und Geistesgegenwart, als sein Herr ihm je zugetraut hätte. Er duckte sich hinter den Schild, sprang mit dem Knüttel gegen den Schützen vor und hieb so rasch um sich, dass einer den Bogen fallenließ und ein anderer mit blutigem Mund rückwärts in den Kot taumelte. Da die Bauern in der Finsternis nicht wußten, ob hinter den beiden noch mehr Räuber stünden, packte sie die Furcht. Sie ließen von \gls{Kuonrat von Steinare} ab, riefen einander in ihrer Sprache zu und flohen über den Hofzaun hinaus in die dunklen Felder.

Als der Sack endlich von seinem Kopf gerissen war, rang \gls{Kuonrat von Steinare} zuerst nach Luft und dann nach Zorn. Statt aber Gott zu danken, dass er nicht mit Stricken gebunden fortgeschleppt worden war, fuhr er den \gls{Knecht} an: \enquote{Du stumpfer Ochse, du hast sie auseinandergejagt. Eine einzige Magd hätte ich gebraucht, und du lässest mir den Hof leerlaufen wie einen aufgescheuchten Hühnerstall.} Der \gls{Knecht} schwieg, wie er fast immer schwieg.

Indessen war die Not der Bauern nicht beim Davonlaufen stehengeblieben. Noch ehe sie den Hof gestürmt hatten, hatten sie den jüngsten Burschen durch einen rückwärtigen Garten hinausgeschickt, auf dass er zur Wache von \gls{Znaim} eile. Dort saßen in jener Nacht Männer bereit, denn auch in der Stadt war das Gerücht von den geraubten Frauen längst angekommen. Also ritten bewaffnete Stadtknechte in kleinen Rotten aus und streiften die Wege zwischen den Weilern, wo \gls{Gruober} und seine Gesellen auf leichte Beute gehofft hatten.

Es währte nicht lange, da stieß einer von \gls{Gruober}s Leuten zu ihnen. Er berichtete, die Wache von \gls{Znaim} habe den alten Spielmann samt zwei Männern auf einem Hohlweg gestellt, als sie eben aus einem anderen Weiler kamen. Einer sei niedergestochen worden, zwei seien in die Finsternis entronnen, doch \gls{Gruober} selbst sei gefangen worden. Man habe ihn, so der Bote schwor, noch in derselben Nacht hinter die Mauern von \gls{Znaim} gebracht.

Ein verständiger Mensch hätte darin ein Zeichen gesehen, die eigene Haut zu retten und sich weit von Stadt, Galgen und Gericht zu entfernen. \gls{Kuonrat von Steinare} jedoch hörte nur, dass ein Mann aus seinem Zuge gefangen worden sei, und empfand darüber weniger Furcht als verletzten Trotz. Denn in seiner verkehrten Art dünkte ihn nicht \gls{Gruober}s Schandtat das Ärgernis, sondern dass fremde Wächter einen fahrenden Mann gegriffen hatten, als gehöre ihnen die nächtliche Straße. \enquote{Nach \gls{Znaim} reiten wir}, sprach er. \enquote{Ich will sehen, ob sie ihn behalten können.} So wandte er noch vor Tagesanbruch das Pferd gegen die Stadt.
